\section{Towards a Global View of Nuclear Structure}
\begin{example}{2}{4-4: J and M Coupling}
\allbold{Using the m scheme, deduce the possible couplings in terms of \(J\) and \(M\) for:}

\allbold{a.) Two identical fermions with \(j = 7/2\) e.g. coupling two \(\text{f}_{7/2}\) neutrons.}
Each particle will have eight possible states (projections), \(\pm \frac{1}{2}, \pm \frac{3}{2}, \pm \frac{5}{2}, \pm \frac{7}{2}\). Since the particles are identical, they cannot occupy the same \(m\) states so \(m_1 \neq m_2\) 



\allbold{b.) Three identical fermions with \(j = 7/2\) and obtained the allowed spins for the \(v=3\) states. Compare your result with information in figure \ref{Tb}.}

\allbold{c.) Two non-identical fermions with \(j = 7/2\) e.g. coupling an \(\text{f}_{7/2}\) proton and \(\text{f}_{7/2}\) neutron.}
    
\end{example}

\begin{example}{3}{4-10: Mirror Nuclei on the \(A=42\) Isobar}
\allbold{Make a comparison of the observed energies in \ensdf for the states in \ce{^{42}Ca}, \ce{^{42}Sc} and \ce{^{42}Ti} below 1600 keV as a side-by-side plot. Noting your answers to 4-4a, and that fact that \ce{^{42}_{20}Ca} and \ce{^{42}_{22}Ti} are a mirror pair, what can you conclude about the behaviour of the two valence nucleons in these \(A=42\) isobars?}
\end{example}

\begin{example}{3}{4-11: Exploring the Structure of \ce{^{207}Pb} and \ce{^{208}Pb} with ENSDF Pages}
\allbold{The following are notes from video exercises exploring different aspects of the structure of \ce{^{207}Pb} intended to encourage familiarity with identifying states adjoining the doubly magic nucleus, \ce{^{208}Pb}.}
\end{example}

\begin{example}{3}{4-6: Classification of the Lowest States in \ce{^{209}Bi}}
\allbold{Using \ensdf and information from the figures below, classify the lowest states in \ce{^{209}Bi} i.e. shell-model particle states, shell-model hole states and core-excited states coupled to shell-model particle states.}
\end{example}

